\chapter[Excerpte zur pastoralen Casuistik]{Excerpte als Beiträge zur pastoralen Casuistik}
\quellseite{26}

\section{Bedenken Dr. Polycarp Leysers}
\subsection[Ob ein Prediger einem Unwürdigen]{Ob ein Prediger wissentlich einem Unwürdigen das heilige Abendmahl reichen und sich mit Christi und Juda’s Exempel behelfen möge?}

Auf vorgelegte Frage eine richtige und dem Glauben ähnliche Antwort zu geben, wird vor allen Dingen erfordert, daß zuvor die zwei Termini unwürdig und wissentlich erklärt werden.

Durch Unwürdige verstehen wir nicht die Schwachgläubigen mit Bucer, sondern die Ungläubigen, die sich nicht selbst prüfen, ohne wahre Reue und Buße, Glauben an Christum und guten Vorsatz, mehr aus Gewohnheit als Andacht sich einstellen und nicht unterscheiden den Leib des Herrn, 1 Cor. 11, 29.

Durch das Wort wissentlich kann beides, ein bestimmtes Wissen und auch eine Meinung, welche sich auf gewisse Vermuthungen gründet, verstanden werden. Denn 1. vielmals einer einer irrigen Meinung in Glaubenssachen oder einer bösen That beschuldigt wird; es sind starke Vermuthungen vorhanden; er hat sich selbst verdächtig gemacht: man kann’s ihm aber zur Genüge nicht überführen. 2. Etliche sind offenbare Sünder, deren Mißhandlungen der ganzen Gemeine kund und offenbar sind. 3. Etliche Sünder sind Gott und dem Prediger allein wissend und bekannt, daß der Prediger entweder die Mißhandlung selbst gesehen, oder wohl von dem Thäter in der Beichte dem Prediger ist offenbaret worden.

Hierauf wird gefragt: Ob ein Prediger wissentlich einem Unwürdigen das Nachtmahl reichen möge? Und ist unterschiedlich zu antworten:

Was erstlich öffentliche unbußfertige Sünder anlanget, welche in ihren Sünden beharren und kein Zeichen wahrer Reue von sich geben, die sind entweder durch die Kirche in den Bann gethan; und alsdann kann ein Prediger ihnen das heilige Nachtmahl nicht reichen, es sei denn, daß sie öffentlich mit der Kirche ausgesühnet sind worden. Da auch ein Prediger, außer eines Todesfalles und wenn er ernstlich Buße gethan hat, einem solchen Gebannten wissentlich das heilige Nachtmahl reichen wollte, hat das gottselige Alterthum ganz eifrig geschlossen, quod sit deponendus, man soll einen solchen Prediger des Dienstes entsetzen. Aber wann die Kirche fahrlässig ist in der Kirchenzucht, muß doch ein Prediger sein Gewissen in Acht nehmen und mag ihnen das heilige Nachtmahl nicht reichen, wie aus den Schriften der Unsrigen zu ersehen.

I. Wegen des Befehls Matth. 7, 6: „Ihr sollt das Heiligthum nicht den Hunden geben und eure Perlen sollt ihr nicht für die Säue werfen.“ Und der Apostel 1 Cor. 5, 11. gebeut den Christen: „So jemand ist, der sich einen Bruder lässet nennen, und ist ein Hurer oder ein Geiziger \&c., mit demselben
\quellseite{27}
sollt ihr auch nicht essen.“ Aus welchem a minori ad majus (vermöge eines Schlusses von dem minder Wichtigen auf das Wichtigere) geschlossen wird: Soll man mit öffentlichen Sündern nicht in gemein essen, vielweniger wird man mit und neben ihnen an dem heiligen Tische des HErrn sich dürfen finden lassen.

II. Wegen des Beispiels. 4 Mos. 9, 6. lesen wir, daß die Unreinen nicht konnten mit den andern Juden das Passah schlachten und essen, sondern sie mußten sich zuvor reinigen und im folgenden Monden das Fest erst halten. Nun hat Christus anstatt der österlichen Mahlzeit das heilige Nachtmahl verordnet, will sich demnach nicht gebühren, daß man, die geistlich unrein sind und in Sünden leben, dazu kommen lasse.

III. Wegen der Gefahr sowohl deßjenigen, der das Abendmahl darreicht, als auch dessen, der dasselbe empfängt. Der Prediger macht sich fremder Sünden theilhaftig, dafür er sich billig hüten sollte: und der Communicant isset und trinket ihm selber das Gericht, 1 Cor. 11, 29. Prediger aber sind von Gott dergestalt zu Richtern gesetzt, daß sie die Gottlosen an Gottes Statt warnen sollen. Wenn sie das nicht thun, sondern den Bösen gleich einem Frommen zum heiligen Nachtmahl zulassen, so stirbt der Gottlose zwar um seiner Sünde willen, aber sein Blut wird aus des Predigers Händen gefordert, Ezech. 3, 16. Damit nun ein Prediger seine Seele errette, so mag er sich in diesem hohen Werck wohl fürsehen und dergleichen Personen nicht zulassen, oder folgen dem Rath des Chrysostomus, welcher Hom. 83. in Matth. also schreibet: „Wenn du ihn aber selbst nicht zurückzuweisen wagst, so magst du es mir sagen, ich werde nicht zulassen, daß dieses geschehe. \gesperrt{Ich will lieber mein Leben hingeben, als daß der Leib des HErrn irgend einem unwürdig gereicht werde, und lieber mein Blut vergießen lassen, als daß ich jenes allerheiligste Blut einem Unwürdigen zulasse}.“\footnote{Chrys. Hom. 83. in Matth.: „Quod si ipse pellere non audes, mihi dicas non permittam ista fieri. Animam potius tradam meam, quam Dominicum corpus alicui dari indigne, et sanguinem meum potius effundi patiar, quam sanctissimum illum sanguinem indigne concedam.“}

Was den andern Haufen anlanget, welche irriger Lehre, als des Calvinismus, oder böser Thaten beschuldigt werden, sind auch starke Anzeigen wider sie vorhanden; sie sind es aber noch nicht geständig und es ist noch nicht über sie ausgeführet: solche Leute kann man nicht einfach abweisen, sondern man soll sie fleißig erinnern, was für große Gefahr ihnen darauf stehe, wenn sie unwürdig von des HErrn Brod essen würden; und sie ermahnen, so lange davon abzustehen, bis sie sich gebührlich reinigen und des Verdachtes entschuldiget werden. Sollte aber solches Vermahnen nichts helfen wollen, sondern sie sich unter die andern Communicanten mit einmengen, weil es heißt: „Aber Verborgenes richtet die Kirche nicht, viel weniger über Unbekanntes“,\footnote{„De occultis non judicat ecclesia, multo minus de incognitis.“} was die gewisse Kenntniß anlangt, so kann ein Prediger ohne Verlegung seines Gewissens ihnen das heilige Nachtmahl reichen, obgleich in künftiger Zeit ein anders wider sie ausgeführt würde.

Denn weil sie getaufte Christen sein, sich als bußfertige Sünder einstellen, so haben sie ein Recht an den gesegneten Tisch und können die Prediger, weil sie nicht Herzenkündiger sind, ihnen solches nicht wehren. Denn die öffentliche Abweisung vom Altar erfordert offenbare und öffentliche Vergehen, welches wir im nachfolgenden Abschnitt ausführlicher machen wollen.
\quellseite{28}
Endlich der dritte Haufe begreift in sich diejenigen, welcher Sünden noch verborgen sind, das ist, welche dem Thäter und Prediger allein bekannt sind. In diesem Fall kann ein Prediger einen solchen, der es bekennt, weder einfach zulassen, noch öffentlich abwehren. Das ist, er soll dem Sünder ernstlich zureden und wegen des Schadens seiner Seele anmahnen, daß er sich enthalte vom heiligen Nachtmahl, bis er rechtschaffene Früchte der Buße leuchten lasse, damit er nicht in Sünden wider das Gewissen sich einstelle und schuldig werde an dem Leibe und Blute des Herrn Christi, auch ihm Mittel und Wege weisen, wie er wahre Buße thun und von der Sünde sich los machen könne, welchem treuen Rath zu folgen er verspricht, wird er billig absolvirt und zugelassen. Würde er aber solcher Zusage nicht nachkommen und die Vermahnung in Wind schlagen und unter andern Communicanten sich einstellen, kann ihn ein Prediger öffentlich nicht abweisen; denn es ist ein verborgenes Vergeben, und ob es wohl dem Prediger offenbar worden, so soll er doch nicht ein Offenbarer des Bekenntnisses sein, welches geschehen würde, wenn er ihn zur Communion nicht zulassen wollte und bei seinen Mitbrüdern ihn dadurch unrichtig machte. Als zum Exempel, wenn eine Dirne wider das sechste Gebot gesündigt hätte, weil sie aber nicht schweres Leibes und niemand die That wissend, behielte sie den Kranz auf, entdeckte aber ihren Fall dem Seelsorger, welcher sie zwar vermahnen soll, den Kranz abzulegen, sintemal sie an Christi Tisch kommen wolle, dem ihre That wissend und sich nicht betrügen läßt, sondern sie gewißlich strafen würde, wäre demnach besser, zeitliche Schande leiden, als ewige gewarten. Wenn sie aber nichts weniger im Kranze sich einstellen würde, kann sie der Prediger nicht öffentlich abweisen, denn jedermann dadurch kund würde, daß sie einer Missethat schuldig, ein Prediger aber seiner Beichtkinder Verräther nicht sein kann noch soll.

Aus diesem Allem steht man auch, wie fern sich ein Prediger mit dem Exempel Christi und Juda’s entschuldigen könne. Denn obwohl etliche in der Meinung seind, daß Judas schon die Gemeinschaft der Apostel verlassen habe, da der Herr das heilige Abendmahl eingesetzt, und damit die Anwendung ganz aufheben; so lassen wir doch gerne zu, daß er das heilige Abendmahl mit den andern Aposteln genossen habe, weil der Evangelist Lucas Cap. 22. solches klarlich bestätigt. Allein es heißt: „Nicht auf Beispielen, sondern auf Gesetzen muß man sehen.“ Sintemal wir nicht wissen, was der Herr Christus für Ursache gehabt habe, darum er auch seinem Verräther diesen theuren Schatz mittheilen wollen, hat uns auch nicht befohlen, solches bei Andern zu thun, sondern vielmehr es verboten, wie oben angeführt worden. Jedoch weil es auch heißt: „Jede Handlung Christi ist unsere Anweisung“;\footnote{Omnis Christi actio est nostra institutio.} ist kein Zweifel, daß wir auch in dieser Handlung des Herrn Christi mit dem Verräther Judas eine gute Unterweisung haben, wie mit seines Gleichen zu gebahren (wohl zu merken, ich sage mit seines Gleichen), denn „Beispiele darf man nicht zu weit und auf Alles ausdehnen“;\footnote{Exempla non sunt nimium et ad omnia extendenda.} sondern man muß in den Grenzen bleiben, weil man sonst ein allgemeines Verbot in der Schrift liest. Und ist von Nöthen, daß wir die Person des Judas wohl betrachten in diesen Puncten. 1. Es war einer aus den Aposteln, er hält sich auch noch zu der Gemeinschaft der Apostel als ein Jünger Jesu, der seiner Lehre öffentlich beipflichtet. 2. Die That war allein Christo bekannt, den übrigen Aposteln verborgen. Denn in der Kirche Christi wußte niemand nichts darum, als
\quellseite{29}
der HErr, und zwar als ein Herzenkündiger; die Hohenpriester wußten es zwar auch, aber sie waren nicht aus den Schäflein Christi. 3. Die That des Verraths war noch nicht zu Werck verrichtet, sondern er hat es im Sinn und den Pharisäern zugesagt und versprochen. 4. Da er darüber zu Rede gesetzt wird, ist er es im Geringsten nicht geständig. 5. Darum wird er nicht bloß alsbald zugelassen, sondern ernstlich ermahnt, Buße zu thun und davon abzustehen. 6. Darauf stellt er sich ein neben den andern Aposteln und wird von dem HErrn zur Communion gelassen. Wo nun in einem Individuum solche Verhältnisse sich finden, achte ich dafür, daß ein Prediger mit gutem Gewissen wissentlich einem Unwürdigen das Nachtmahl reichen könne und für Gott entschuldiget sei. Und solches desto mehr, weil ein Prediger über die Buße, welche folgen oder ausbleiben wird, ungewiß ist und dem Erbieten der Menschen trauen muß, das Uebrige aber Gott befehlen. Er reicht es ihm aber nicht als einem Unwürdigen, sondern als einem, der zum Schein Unschuld vorgibt.

So geht nun dies Exempel Juda nicht an: 1. die Ketzer, die sich von dem Bekenntniß unserer Kirche öffentlich absondern. 2. Auch nicht notorische Sünder, deren Sünde der Kirche oder doch den Meisten darinnen wissend und bekannt sind. 3. Auch nicht diejenigen, welche eine böse That völlig zu Wercke gerichtet haben. 4. Auch nicht, die ihrer Mißhandlung keine Scheu tragen, öffentlich es gestehen und wohl dazu sich rühmen. Solche soll man nicht allein strafen, sondern auch gar abweisen. Sondern es läßt sich dieses Exempel nur auf die ziehen, die 1. im äußerlichen und öffentlichen Bekenntniß sich zu der Gemeine Christi halten. 2. Vor der Welt ein ehrbares Leben führen und nicht böser Thaten berüchtiget seind. 3. Welche mit Bösem schwanger gehen, dem Nächsten dräuen oder sich Schaden zu thun rühmen. 4. Wenn sie ernstlich zu Rede gesetzt werden, es verneinen und nicht geständig sein wollen. 5. Wenn solche nach ernster Vermahnung und Erinnerung, daß sie vom Bösen abstehen und wahre Buße thun sollen, 6. sich neben andern Christen zur Communion einstellen.

Obwohl ein Prediger an wahrer Bekehrung und Besserung zweifelt, auch aus den Werken ein Widriges geschlossen möchte werden; so soll doch ein Prediger seines Oberherrn Christ Sanftmuth und Geduld folgen, solche zulassen und das Uebrige Gott, dem Herzenkündiger, befehlen. Wer aber dieses Exempel weiter ausdehnen will, der nehme wohl in Acht das allgemeine Verbot Christi und die Analogie des Glaubens, damit er nicht auf einen sandigen Boden baue und in den Versuchungen mit solcher Entschuldigung sein Gewissen mehr beunruhige als aufrichte. Welches ich dem Herrn auf vorgelegte Frage kürzlich antworten wollen, von Herzen wünschend, daß Gott mit seinem Heiligen Geist alle Communicanten wolle regieren, damit sie sich recht prüfen und würdige Gäste am Tische des HErrn sein mögen zu seines Namens Lob und Preis und der Communicanten ewigen Seligkeit, und das wolle er thun um Jesu Christi willen. Amen.

(Siehe: Thesauri Consiliorum Vol. I. Part. 2. In Druck gegeben durch M. Georgium Dedekennum, fol. 351–354.)
